%Tutti gli usepackage vanno qui

\usepackage{geometry}
\usepackage[italian]{babel}
\usepackage[utf8]{inputenc}
\usepackage[T1]{fontenc}
\usepackage{tabularx}
\usepackage{longtable}
\usepackage{hyperref}
\hypersetup{
	colorlinks=true,
	linkcolor=black,
	filecolor=magenta,      
	urlcolor=blue,
}
\usepackage[bottom]{footmisc}
\usepackage{fancyhdr}
\usepackage{titlesec}
\setcounter{secnumdepth}{4}
\usepackage{amsmath, amssymb}
%\usepackage{caption}
\usepackage{graphicx}
%\usepackage{float}
\usepackage{layouts}
\usepackage{float}


%Comandi di impaginazione uguale per tutti i documenti
\pagestyle{fancy}
\lhead{\includegraphics[scale=0.07]{res/images/logo8_crop.png}}
%Titolo del documento
\rhead{\doctitle{}}
%\lfoot{\rightmark}
\rfoot{\thepage}
\cfoot{}
\setlength{\headheight}{35pt}
\setcounter{tocdepth}{4}
\setcounter{secnumdepth}{4}
\renewcommand{\footrulewidth}{0.4pt}

% multirow per tabelle
\usepackage{multirow}

% Permette tabelle su più pagine
\usepackage{longtable}

% colore di sfondo per le celle
\usepackage[table]{xcolor}

%COMANDI TABELLE
\newcommand{\rowcolorhead}{\rowcolor[HTML]{56A5EC}} %intestazione 
\newcommand{\rowcolorlight}{\rowcolor[HTML]{A0CFEC}} %righe chiare/dispari
\newcommand{\rowcolordark}{\rowcolor[HTML]{EBDDE2}} %righe scure/pari
\newcommand{\colorhead}{\color[HTML]{FFFFFF}} %testo intestazione
\newcommand{\colorbody}{\color[HTML]{000000}} %testo righe
\newcommand{\tabitem}{~~\rlap{\textbullet}~~}